% --- SLIDE 2: QUANTISATION, THE MAPPING ---
\begin{frame}{Quantisation: Fewer Bits per Number}

    \begin{columns}[T]
        \begin{column}{0.5\textwidth}
            \begin{itemize}
                \small
                \item Map each real value $x$ to an integer $x_q$ with a \textbf{scale} $s$ and a
                      \textbf{zero-point} $z$:
                      \[
                          x_q = \mathrm{clip}\big(\mathrm{round}(x/s + z),\, q_{\min}, q_{\max}\big)
                      \]
                      and dequantising gives back $x \approx s\,(x_q - z)$.
                \item \textbf{INT8} has 256 levels, so $s$ comes from the observed range of $x$.
                      \textbf{Symmetric} quantisation fixes $z = 0$.
                \item \textbf{Rounding and clipping lose information}, and dequantising does not get it back.
                      One outlier stretches the range and wastes levels, which is why scales are often kept
                      \textbf{per channel} rather than per tensor.
            \end{itemize}
        \end{column}
        \begin{column}{0.48\textwidth}
            \centering
            % Real values (orange) map to integer levels (green). Redrawn after Wang & Spindler,
            % NVIDIA Technical Blog (2025), Figure 2: the original is too thin to read at this size.
            \resizebox{\linewidth}{!}{%
            \begin{tikzpicture}[>=Stealth, font=\small]
                \node[anchor=west, font=\small\bfseries] at (0,3.55) {Affine (asymmetric)};
                \draw[<->, thick, orange!85!black] (0,3.0) -- (6.2,3.0);
                \draw[thick, green!45!black] (0.3,2.0) -- (5.9,2.0);
                \foreach \t/\b/\tl/\bl in {1.0/0.3/x_{\min}/q_{\min}, 2.2/1.9/0/z, 3.8/3.9/x/x_q, 5.2/5.9/x_{\max}/q_{\max}} {
                    \fill[orange!85!black] (\t,3.0) circle (1.8pt) node[above, black] {$\tl$};
                    \fill[green!45!black] (\b,2.0) circle (1.8pt) node[below, black] {$\bl$};
                    \draw[->, gray] (\t,2.92) -- (\b,2.09);
                }
                \node[anchor=west, font=\small\bfseries] at (0,1.15) {Symmetric};
                \draw[<->, thick, orange!85!black] (0,0.6) -- (6.2,0.6);
                \draw[thick, green!45!black] (0.3,-0.4) -- (5.9,-0.4);
                \foreach \t/\b/\tl/\bl in {1.0/0.3/-x_{\max}/q_{\min}, 3.1/3.1/0/0, 4.2/4.6/x/x_q, 5.2/5.9/x_{\max}/q_{\max}} {
                    \fill[orange!85!black] (\t,0.6) circle (1.8pt) node[above, black] {$\tl$};
                    \fill[green!45!black] (\b,-0.4) circle (1.8pt) node[below, black] {$\bl$};
                    \draw[->, gray] (\t,0.52) -- (\b,-0.31);
                }
            \end{tikzpicture}}\\[0.3em]
            {\scriptsize Real values (top) to integer levels (bottom). After Wang \& Spindler, NVIDIA
            Technical Blog (2025), Figure~2.}
        \end{column}
    \end{columns}

\end{frame}

% --- SLIDE 3: QUANTISATION, WHICH METHOD ---
\begin{frame}{Quantisation: When and How}

    \centering
    \renewcommand{\arraystretch}{1.3}
    {\scriptsize
    \begin{tabular}{p{3.0cm}|p{5.0cm}|p{3.6cm}}
        \toprule
        \textbf{Method} & \textbf{How it works} & \textbf{Use it when} \\
        \midrule
        Post-training, \textbf{dynamic}
            & Weights are quantised ahead of time; activation scales are computed on the fly for each input
            & A first try: needs no data. Linear- and LSTM-heavy models on CPU \\
        Post-training, \textbf{static}
            & A \textbf{calibration set} of representative inputs fixes the activation scales once
            & Fastest at run time; CNNs and edge devices \\
        \textbf{Quantisation-aware training}
            & Training simulates the rounding (\emph{fake quantisation}, with gradients passed straight
              through it), so the weights adapt to it
            & Post-training loses too much accuracy, e.g.\ at 4 bits or on small models \\
        \bottomrule
    \end{tabular}}

    \vspace{0.6em}
    \begin{itemize}
        \item \textbf{The calibration set is training data for the quantiser.} Version it with the model, and
              check it still looks like production when the inputs drift.
        \item \textbf{Tools:} \textbf{torchao} for PyTorch (\texttt{torch.ao.quantization} is deprecated),
              \textbf{ONNX Runtime} quantisation for exported models, \textbf{TensorRT} or
              \textbf{OpenVINO} for NVIDIA or Intel targets.
    \end{itemize}

\end{frame}
